% problem_id: S001-lemma-vanishing-coh-almost-projective
% source_id: S001 (Stacks Project)
% source_file: local-cohomology.tex
% statement_locator: local-cohomology.tex lines 4927--4954
% proof_locator: local-cohomology.tex lines 4956--5073
% env: lemma
% pairing: tight (verified)
% status: complete (statement + author proof verbatim + reason + locators)
%
\begin{lemma}
\label{lemma-vanishing-coh-almost-projective}
With $q_0 = q(S)$ and $d = d(S)$ as above, suppose we have
integers $n \geq c \geq 0$, an $(A, n, c)$-module $M$,
an index $i \in \{0, 1, \ldots, d\}$, and an integer $q$.
Then we distinguish the following cases
\begin{enumerate}
\item In the case $i = d \geq 1$ and $q \geq q_0$ we have
$H^d(X, p^*M(q)) = 0$.
\item In the case $i = d - 1 \geq 1$ and $q \geq q_0$ we have
$H^{d - 1}(X, p^*M(q)) = 0$.
\item In the case $d - 1 > i > 0$ and $q \geq q_0 + (d - 1 - i)c$
the map
$H^i(X, p^*M(q)) \to H^i(X, p^*M(q - (d - 1 - i)c))$
is zero.
\item In the case $i = 0$, $d \in \{0, 1\}$, and $q \geq q_0$, there
is a surjection
$$
I^qM \longrightarrow H^0(X, p^*M(q))
$$
\item In the case $i = 0$, $d > 1$, and $q \geq q_0 + (d - 1)c$ the map
$$
H^0(X, p^*M(q)) \to H^0(X, p^*M(q - (d - 1)c))
$$
has image contained in the image of the canonical map
$I^{q - (d - 1)c}M \to H^0(X, p^*M(q - (d - 1)c))$.
\end{enumerate}
\end{lemma}

\begin{proof}
Let $M$ be an $(A, n, c)$-module. Choose a short exact sequence
$$
0 \to K \to (A/I^n)^{\oplus r} \to M \to 0
$$
We will use below that $K$ is an $(A, n, c)$-module, see More on Algebra,
Lemma \ref{more-algebra-lemma-ses-near-projective}.
Consider the corresponding exact sequence
$$
p^*K \to (\mathcal{O}_{Y_n})^{\oplus r} \to p^*M \to 0
$$
We split this into short exact sequences
$$
0 \to \mathcal{F} \to p^*K \to \mathcal{G} \to 0
\quad\text{and}\quad
0 \to \mathcal{G} \to (\mathcal{O}_{Y_n})^{\oplus r} \to p^*M \to 0
$$
By Lemma \ref{lemma-almost-exactness} the coherent module $\mathcal{F}$
is scheme theoretically supported on $Y_c$.

\medskip\noindent
Proof of (1). Assume $d > 0$. We have to prove
$H^d(X, p^*M(q)) = 0$ for $q \geq q_0$.
By the vanishing of the cohomology of twists of $\mathcal{G}$ in degrees $> d$
and the long exact cohomology sequence associated to the second
short exact sequence above, it suffices to prove that
$H^d(X, \mathcal{O}_{Y_n}(q)) = 0$.
This is true by Lemma \ref{lemma-bound-q-and-d}.

\medskip\noindent
Proof of (2). Assume $d > 1$. We have to prove
$H^{d - 1}(X, p^*M(q)) = 0$ for $q \geq q_0$.
Arguing as in the previous paragraph, we see that it suffices
to show that $H^d(X, \mathcal{G}(q)) = 0$. Using the first
short exact sequence and the vanishing of the cohomology
of twists of $\mathcal{F}$ in degrees $> d$ we see that it suffices
to show $H^d(X, p^*K(q))$ is zero which is
true by (1) and the fact that $K$ is an $(A, n, c)$-module (see above).

\medskip\noindent
Proof of (3). Let $0 < i < d - 1$ and assume the statement holds for $i + 1$
except in the case $i = d - 2$ we have statement (2).
Using the long exact sequence of cohomology associated to the second
short exact sequence above we find an injection
$$
H^i(X, p^*M(q - (d - 1 - i)c)) \subset
H^{i + 1}(X, \mathcal{G}(q - (d - 1 - i)c))
$$
as $q - (d - 1 - i)c \geq q_0$ gives the vanishing of
$H^i(X, \mathcal{O}_{Y_n}(q - (d - 1 - i)c))$
(see above). Thus it suffices to show that the map
$H^{i + 1}(X, \mathcal{G}(q)) \to H^{i + 1}(X, \mathcal{G}(q - (d - 1 - i)c))$
is zero. To study this, we consider the maps of exact sequences
$$
\xymatrix{
H^{i + 1}(X, p^*K(q)) \ar[r] \ar[d] &
H^{i + 1}(X, \mathcal{G}(q)) \ar[r] \ar[d] \ar@{..>}[ld] &
H^{i + 2}(X, \mathcal{F}(q)) \ar[d] \\
H^{i + 1}(X, p^*K(q - c)) \ar[r] \ar[d] &
H^{i + 1}(X, \mathcal{G}(q - c)) \ar[r] \ar[d] &
H^{i + 2}(X, \mathcal{F}(q - c)) \\
H^{i + 1}(X, p^*K(q - (d - 1 - i)c)) \ar[r] &
H^{i + 1}(X, \mathcal{G}(q - (d - 1 - i)c))
}
$$
Since $\mathcal{F}$ is scheme theoretically supported on $Y_c$
we see that the canonical map
$\mathcal{G}(q) \to \mathcal{G}(q - c)$ factors through
$p^*K(q - c)$ by Lemma \ref{lemma-annihilated}.
This gives the dotted arrow in the diagram. (In fact, for the proof it
suffices to observe that the vertical arrow on the extreme right is
zero in order to get the dotted arrow as a map of sets.)
Thus it suffices to show that
$H^{i + 1}(X, p^*K(q - c)) \to
H^{i + 1}(X, p^*K(q - (d - 1 - i)c))$
is zero. If $i = d - 2$, then the source of this arrow
is zero by (2) as $q - c \geq q_0$ and $K$ is an $(A, n, c)$-module.
If $i < d - 2$, then as $K$ is an $(A, n, c)$-module, we get from the
induction hypothesis that the map is indeed zero
since $q - c - (q - (d - 1 - i)c) = (d - 2 - i)c = (d - 1 - (i + 1))c$
and since $q - c \geq q_0 + (d - 1 - (i + 1))c$.
In this way we conclude the proof of (3).

\medskip\noindent
Proof of (4). Assume $d \in \{0, 1\}$ and $q \geq q_0$.
Then the first short exact sequence gives a surjection
$H^1(X, p^*K(q)) \to H^1(X, \mathcal{G}(q))$
and the source of this arrow is zero by case (1). Hence
for all $q \in \mathbf{Z}$ we see that the map
$$
H^0(X, (\mathcal{O}_{Y_n})^{\oplus r}(q))
\longrightarrow
H^0(X, p^*M(q))
$$
is surjective. For $q \geq q_0$ the
source is equal to $(I^q/I^{q + n})^{\oplus r}$ by
Lemma \ref{lemma-bound-q-and-d} and this easily proves the statement.

\medskip\noindent
Proof of (5). Assume $d > 1$. Arguing as in the proof of (4) we see that
it suffices to show that the image of
$$
H^0(X, p^*M(q))
\longrightarrow
H^0(X, p^*M(q - (d - 1)c))
$$
is contained in the image of
$$
H^0(X, (\mathcal{O}_{Y_n})^{\oplus r}(q - (d - 1)c))
\longrightarrow
H^0(X, p^*M(q - (d - 1)c))
$$
To show the inclusion above, it suffices to show that for
$\sigma \in H^0(X, p^*M(q))$ with boundary
$\xi \in H^1(X, \mathcal{G}(q))$ the image of $\xi$ in
$H^1(X, \mathcal{G}(q - (d - 1)c))$ is zero. This follows by the
exact same arguments as in the proof of (3).
\end{proof}
